\begin{casebox}
\textbf{知识点小案例：响应丢失会把“执行一次”变成协议问题。}
特例是 worker 已经写出 block，但 reply 在网络中丢失；controller 超时后启动 attempt 1，attempt 1 也写出 block 并提交。稍后 attempt 0 的迟到 reply 到达。若消息里没有 task id、attempt id、epoch、deadline 和输出身份，controller 无法判断该接受还是拒绝。由此推出规律：分布式系统不承诺“调用一次执行一次”，它承诺在身份、重试和提交表约束下做出可恢复决定。工程上每个 reply 只能进入 classify，再由 commit 规则决定是否写 manifest。
\end{casebox}
